% !TeX program = lualatex
% !TeX root = predavanje06.tex
% ============================================================
% Zajednički preamble za sva predavanja RWA kolegija
% Uključiti s: % ============================================================
% Zajednički preamble za sva predavanja RWA kolegija
% Uključiti s: % ============================================================
% Zajednički preamble za sva predavanja RWA kolegija
% Uključiti s: \input{../common/preamble.tex}
% ============================================================

\documentclass[aspectratio=149,11pt]{beamer}
\usetheme{CambridgeUS}

\usepackage{fontspec}              % radi samo sa XeLaTeX/LuaLaTeX
\setsansfont{Latin Modern Sans}    % sans-serif font
\setmonofont{Latin Modern Mono}    % monospace font za kod
\usepackage[croatian]{babel}
\usepackage{csquotes}              % \enquote{} → hrvatski navodnici „..."
\usepackage{lmodern}
\usepackage{booktabs}
\usepackage{listings}
\usepackage{xcolor}
\usepackage{hyperref}

% ---------- Autor i logo ----------
\author{mag.ing. Josip Torić (SIT UNIZD)}
\titlegraphic{\includegraphics[height=2.5cm]{../common/SIT_logo_tekst1.pdf}}

% ---------- Listings ----------
\definecolor{codebg}{RGB}{248,248,248}
\definecolor{codeframeclr}{RGB}{220,220,220}
\lstdefinelanguage{bash}{
  morekeywords={sudo,apt,brew,pip,pipx,uv,winget,python,git,flask,source,deactivate,code,gunicorn,docker,compose,heroku,render,railway,node,npm,pnpm,yarn,uvicorn,alembic,pytest,ruff,black,precommit},
  sensitive=false,
  morecomment=[l]{\#}
}
\lstdefinelanguage{python}{
  morekeywords={from,import,def,class,return,if,elif,else,for,in,while,try,except,raise,with,as,not,and,or,is,lambda,yield,True,False,None,await,async},
  sensitive=true,
  morecomment=[l]{\#},
  morestring=[b]',
  morestring=[b]"
}
\lstdefinelanguage{sql}{
  morekeywords={SELECT,FROM,WHERE,INSERT,INTO,VALUES,UPDATE,SET,DELETE,CREATE,TABLE,PRIMARY,KEY,FOREIGN,REFERENCES,UNIQUE,NOT,NULL,INDEX,ON,ALTER,DROP,CONSTRAINT},
  sensitive=false,
  morecomment=[l]{--}
}
\lstdefinelanguage{json}{
  morecomment=[l]{//},
  morestring=[b]",
  showstringspaces=false
}
\lstdefinelanguage{typescript}{
  morekeywords={import,export,from,const,let,var,function,return,if,else,for,of,in,while,switch,case,break,async,await,try,catch,finally,throw,new,class,this,extends,implements,interface,type,enum,readonly,public,private,protected,static,get,set,true,false,null,undefined,void,string,number,boolean,any,unknown,never,typeof,keyof,as},
  sensitive=true,
  morecomment=[l]{//},
  morecomment=[s]{/*}{*/},
  morestring=[b]',
  morestring=[b]",
  morestring=[b]`
}
\lstalias{javascript}{typescript}
\lstset{
  backgroundcolor=\color{codebg},
  basicstyle=\ttfamily\footnotesize,
  frame=single,
  rulecolor=\color{codeframeclr},
  frameround=tttt,
  breaklines=true,
  tabsize=2,
  showstringspaces=false,
  upquote=true,
  columns=fullflexible
}

% ---------- Helper macros ----------
\newcommand{\codefile}[1]{\texttt{#1}}
\newcommand{\kbd}[1]{\texttt{#1}}

% ---------- Code frame environment ----------
% Automatski dodaje fragile i shrink=5, koristi se umjesto \begin{frame}
% Primjer: \begin{codeframe}{Naslov} ... \end{codeframe}
% Primjer s opcijama: \begin{codeframe}[shrink=15]{Naslov} ... \end{codeframe}
\newenvironment{codeframe}[2][]
  {\begin{frame}[fragile,shrink=5,#1]{#2}}
  {\end{frame}}

% ============================================================

\documentclass[aspectratio=149,11pt]{beamer}
\usetheme{CambridgeUS}

\usepackage{fontspec}              % radi samo sa XeLaTeX/LuaLaTeX
\setsansfont{Latin Modern Sans}    % sans-serif font
\setmonofont{Latin Modern Mono}    % monospace font za kod
\usepackage[croatian]{babel}
\usepackage{csquotes}              % \enquote{} → hrvatski navodnici „..."
\usepackage{lmodern}
\usepackage{booktabs}
\usepackage{listings}
\usepackage{xcolor}
\usepackage{hyperref}

% ---------- Autor i logo ----------
\author{mag.ing. Josip Torić (SIT UNIZD)}
\titlegraphic{\includegraphics[height=2.5cm]{../common/SIT_logo_tekst1.pdf}}

% ---------- Listings ----------
\definecolor{codebg}{RGB}{248,248,248}
\definecolor{codeframeclr}{RGB}{220,220,220}
\lstdefinelanguage{bash}{
  morekeywords={sudo,apt,brew,pip,pipx,uv,winget,python,git,flask,source,deactivate,code,gunicorn,docker,compose,heroku,render,railway,node,npm,pnpm,yarn,uvicorn,alembic,pytest,ruff,black,precommit},
  sensitive=false,
  morecomment=[l]{\#}
}
\lstdefinelanguage{python}{
  morekeywords={from,import,def,class,return,if,elif,else,for,in,while,try,except,raise,with,as,not,and,or,is,lambda,yield,True,False,None,await,async},
  sensitive=true,
  morecomment=[l]{\#},
  morestring=[b]',
  morestring=[b]"
}
\lstdefinelanguage{sql}{
  morekeywords={SELECT,FROM,WHERE,INSERT,INTO,VALUES,UPDATE,SET,DELETE,CREATE,TABLE,PRIMARY,KEY,FOREIGN,REFERENCES,UNIQUE,NOT,NULL,INDEX,ON,ALTER,DROP,CONSTRAINT},
  sensitive=false,
  morecomment=[l]{--}
}
\lstdefinelanguage{json}{
  morecomment=[l]{//},
  morestring=[b]",
  showstringspaces=false
}
\lstdefinelanguage{typescript}{
  morekeywords={import,export,from,const,let,var,function,return,if,else,for,of,in,while,switch,case,break,async,await,try,catch,finally,throw,new,class,this,extends,implements,interface,type,enum,readonly,public,private,protected,static,get,set,true,false,null,undefined,void,string,number,boolean,any,unknown,never,typeof,keyof,as},
  sensitive=true,
  morecomment=[l]{//},
  morecomment=[s]{/*}{*/},
  morestring=[b]',
  morestring=[b]",
  morestring=[b]`
}
\lstalias{javascript}{typescript}
\lstset{
  backgroundcolor=\color{codebg},
  basicstyle=\ttfamily\footnotesize,
  frame=single,
  rulecolor=\color{codeframeclr},
  frameround=tttt,
  breaklines=true,
  tabsize=2,
  showstringspaces=false,
  upquote=true,
  columns=fullflexible
}

% ---------- Helper macros ----------
\newcommand{\codefile}[1]{\texttt{#1}}
\newcommand{\kbd}[1]{\texttt{#1}}

% ---------- Code frame environment ----------
% Automatski dodaje fragile i shrink=5, koristi se umjesto \begin{frame}
% Primjer: \begin{codeframe}{Naslov} ... \end{codeframe}
% Primjer s opcijama: \begin{codeframe}[shrink=15]{Naslov} ... \end{codeframe}
\newenvironment{codeframe}[2][]
  {\begin{frame}[fragile,shrink=5,#1]{#2}}
  {\end{frame}}

% ============================================================

\documentclass[aspectratio=149,11pt]{beamer}
\usetheme{CambridgeUS}

\usepackage{fontspec}              % radi samo sa XeLaTeX/LuaLaTeX
\setsansfont{Latin Modern Sans}    % sans-serif font
\setmonofont{Latin Modern Mono}    % monospace font za kod
\usepackage[croatian]{babel}
\usepackage{csquotes}              % \enquote{} → hrvatski navodnici „..."
\usepackage{lmodern}
\usepackage{booktabs}
\usepackage{listings}
\usepackage{xcolor}
\usepackage{hyperref}

% ---------- Autor i logo ----------
\author{mag.ing. Josip Torić (SIT UNIZD)}
\titlegraphic{\includegraphics[height=2.5cm]{../common/SIT_logo_tekst1.pdf}}

% ---------- Listings ----------
\definecolor{codebg}{RGB}{248,248,248}
\definecolor{codeframeclr}{RGB}{220,220,220}
\lstdefinelanguage{bash}{
  morekeywords={sudo,apt,brew,pip,pipx,uv,winget,python,git,flask,source,deactivate,code,gunicorn,docker,compose,heroku,render,railway,node,npm,pnpm,yarn,uvicorn,alembic,pytest,ruff,black,precommit},
  sensitive=false,
  morecomment=[l]{\#}
}
\lstdefinelanguage{python}{
  morekeywords={from,import,def,class,return,if,elif,else,for,in,while,try,except,raise,with,as,not,and,or,is,lambda,yield,True,False,None,await,async},
  sensitive=true,
  morecomment=[l]{\#},
  morestring=[b]',
  morestring=[b]"
}
\lstdefinelanguage{sql}{
  morekeywords={SELECT,FROM,WHERE,INSERT,INTO,VALUES,UPDATE,SET,DELETE,CREATE,TABLE,PRIMARY,KEY,FOREIGN,REFERENCES,UNIQUE,NOT,NULL,INDEX,ON,ALTER,DROP,CONSTRAINT},
  sensitive=false,
  morecomment=[l]{--}
}
\lstdefinelanguage{json}{
  morecomment=[l]{//},
  morestring=[b]",
  showstringspaces=false
}
\lstdefinelanguage{typescript}{
  morekeywords={import,export,from,const,let,var,function,return,if,else,for,of,in,while,switch,case,break,async,await,try,catch,finally,throw,new,class,this,extends,implements,interface,type,enum,readonly,public,private,protected,static,get,set,true,false,null,undefined,void,string,number,boolean,any,unknown,never,typeof,keyof,as},
  sensitive=true,
  morecomment=[l]{//},
  morecomment=[s]{/*}{*/},
  morestring=[b]',
  morestring=[b]",
  morestring=[b]`
}
\lstalias{javascript}{typescript}
\lstset{
  backgroundcolor=\color{codebg},
  basicstyle=\ttfamily\footnotesize,
  frame=single,
  rulecolor=\color{codeframeclr},
  frameround=tttt,
  breaklines=true,
  tabsize=2,
  showstringspaces=false,
  upquote=true,
  columns=fullflexible
}

% ---------- Helper macros ----------
\newcommand{\codefile}[1]{\texttt{#1}}
\newcommand{\kbd}[1]{\texttt{#1}}

% ---------- Code frame environment ----------
% Automatski dodaje fragile i shrink=5, koristi se umjesto \begin{frame}
% Primjer: \begin{codeframe}{Naslov} ... \end{codeframe}
% Primjer s opcijama: \begin{codeframe}[shrink=15]{Naslov} ... \end{codeframe}
\newenvironment{codeframe}[2][]
  {\begin{frame}[fragile,shrink=5,#1]{#2}}
  {\end{frame}}


\title[Programiranje za web]{Programiranje za web - Predavanje 6 \\ Poslovna pravila i workflow u backendu}
\date{09.11.2026.}

\begin{document}

\frame{\titlepage}

% ------------------------------------------------------------
\begin{frame}{Plan predavanja}
\begin{enumerate}
  \item Što su poslovna pravila i zašto ne žive u routeru
  \item Rokovi nominacija: prelim i final deadline
  \item Faze natjecanja i što klub smije u kojoj fazi
  \item Implementacija u service layeru
  \item Domain errori za prekršena pravila
  \item Testiranje poslovne logike
  \item Priprema za obranu 1. dijela
\end{enumerate}

\vspace{0.6em}
Na kraju predavanja: backend \textbf{odbija} operacije koje krše poslovna pravila, s jasnim error porukama.
\end{frame}

% ------------------------------------------------------------
\begin{frame}{Ishodi učenja}
Nakon ovog predavanja student može:
\begin{itemize}
  \item definirati i implementirati vremenski uvjetovana poslovna pravila
  \item objasniti razliku između validacije inputa i poslovnog pravila
  \item strukturirati service metodu koja provjerava fazu natjecanja
  \item napisati test koji verificira ponašanje prije i nakon roka
  \item koristiti \texttt{AppError} za domenske greške s odgovarajućim status codeom
\end{itemize}
\end{frame}

% ------------------------------------------------------------
\begin{frame}{Kontekst: što imamo, što danas zatvaramo}
\begin{columns}[T]
\column{0.45\textwidth}
\textbf{Implementirano (P1-P5)}
\begin{itemize}
  \item Svi modeli i CRUD
  \item Auth (JWT)
  \item Autorizacija (role/ownership)
  \item Validacija i paginacija
\end{itemize}

\column{0.45\textwidth}
\textbf{Danas}
\begin{itemize}
  \item Rokovi i faze
  \item Pravila: tko smije što i kad
  \item Testovi za pravila
  \item Zaokruživanje za obranu
\end{itemize}
\end{columns}
\end{frame}

% ------------------------------------------------------------
\begin{frame}{Poslovna pravila vs validacija inputa}
\begin{tabular}{@{}lp{5cm}p{5cm}@{}}
  \toprule
  & \textbf{Validacija inputa} & \textbf{Poslovno pravilo} \\
  \midrule
  Gdje & Pydantic schema & Service layer \\
  Primjer & \texttt{category} mora biti iz IPF liste & klub ne smije prijaviti liftera nakon roka \\
  Status & 422 & 400/403/409 \\
  Ovisi o & formatu podataka & stanju sustava \\
  \bottomrule
\end{tabular}

\vspace{0.6em}
Pydantic provjerava \textbf{oblik}. Service provjerava \textbf{smisao}.
\end{frame}

% ------------------------------------------------------------
\begin{frame}{Faze natjecanja (ponovljeno s P0)}
\begin{tabular}{@{}lp{9cm}@{}}
  \toprule
  \textbf{Faza} & \textbf{Što klub smije} \\
  \midrule
  Do \textbf{prelim} roka & Dodavati nove prijave, uređivati podatke, odjaviti natjecatelja \\
  \midrule
  Između \textbf{prelim} i \textbf{final} & Odjaviti natjecatelja ili promijeniti kategoriju. \textbf{Ne smije} dodavati nove prijave \\
  \midrule
  Nakon \textbf{final} roka & Nema izmjena \\
  \bottomrule
\end{tabular}

\vspace{0.6em}
Ovo su pravila koja backend mora enforcati. Frontend može sakriti gumbe, ali backend \textbf{mora odbiti}.
\end{frame}

% ------------------------------------------------------------
\begin{frame}{State machine: faze kao konačni automat}
Faze natjecanja su \textbf{konačni automat} (finite state machine):

\vspace{0.6em}
\begin{center}
OPEN $\xrightarrow{\text{prelim\_deadline}}$ PRELIM\_PASSED $\xrightarrow{\text{final\_deadline}}$ CLOSED
\end{center}

\vspace{0.6em}
Svojstva:
\begin{itemize}
  \item Prijelazi su \textbf{jednosmjerni} - ne možemo se vratiti u OPEN
  \item Prijelaz ovisi o \textbf{vremenu}, ne o korisničkoj akciji
  \item Svaka faza ima jasno definirane \textbf{dozvoljene operacije}
  \item Nema dvosmislenosti - u svakom trenutku faza je točno jedna
\end{itemize}

\vspace{0.4em}
Kad pravila postanu složenija (npr. 5+ stanja, korisničke tranzicije), razmislite o biblioteci za state machine.
\end{frame}

% ------------------------------------------------------------
\begin{frame}{Soft delete vs hard delete}
\begin{tabular}{@{}lp{5cm}p{5cm}@{}}
  \toprule
  & \textbf{Hard delete} (\texttt{DELETE}) & \textbf{Soft delete} (status) \\
  \midrule
  Što se dogodi & Red nestaje iz baze & Red ostaje, \texttt{status} se mijenja \\
  Audit trail & Nema (osim logova) & Potpun - vidimo tko je i kad odjavljen \\
  Undo & Nemoguće & Moguće (vrati status na \texttt{"active"}) \\
  Query složenost & Jednostavni & Treba filtrirati po statusu \\
  \bottomrule
\end{tabular}

\vspace{0.6em}
Naš projekt koristi soft delete za registracije (\texttt{status = "withdrawn"}) jer:
\begin{itemize}
  \item Organizator treba znati da je netko bio prijavljen pa se odjavio
  \item Sprječava ponovnu prijavu istog liftera (unique constraint i dalje vrijedi)
\end{itemize}
\end{frame}

% ------------------------------------------------------------
\begin{frame}{Idempotentnost: što ako korisnik klikne dva puta?}
Operacija je \textbf{idempotentna} ako ponovljeni poziv daje isti rezultat.

\vspace{0.6em}
\begin{tabular}{@{}lll@{}}
  \toprule
  \textbf{Operacija} & \textbf{Idempotentna?} & \textbf{Naša implementacija} \\
  \midrule
  GET lista & Da & Uvijek isti rezultat \\
  POST registracija & Ne & Drugi poziv: 409 (duplikat) \\
  POST withdraw & Ne & Drugi poziv: 400 (\texttt{already\_withdrawn}) \\
  PATCH kategorija & Da & Isti rezultat koliko god puta \\
  \bottomrule
\end{tabular}

\vspace{0.6em}
Za ne-idempotentne operacije frontend treba spriječiti duple klikove - gumb je disabled dok request traje.
Backend ih odbija s jasnim error codeom - frontend zna da nije greška, nego dupli poziv.
\end{frame}

% ------------------------------------------------------------
\begin{frame}{Određivanje faze: pomoćna funkcija}
Fazu određujemo usporedbom trenutnog vremena s rokovima natjecanja.

\vspace{0.6em}
Tri moguća rezultata:
\begin{enumerate}
  \item \texttt{OPEN} - prije prelim roka
  \item \texttt{PRELIM\_PASSED} - između prelim i final roka
  \item \texttt{CLOSED} - nakon final roka
\end{enumerate}

\vspace{0.6em}
Ova funkcija ne ovisi o korisniku ni o requestu - samo o natjecanju i satu.
\end{frame}

% ------------------------------------------------------------
\begin{frame}[fragile,shrink=5]{Implementacija: faza natjecanja}
\begin{lstlisting}[language=python]
from enum import Enum
from datetime import datetime, timezone

class Phase(str, Enum):
    OPEN = "open"
    PRELIM_PASSED = "prelim_passed"
    CLOSED = "closed"

def get_competition_phase(competition: Competition) -> Phase:
    now = datetime.now(timezone.utc)
    if now < competition.prelim_deadline:
        return Phase.OPEN
    if now < competition.final_deadline:
        return Phase.PRELIM_PASSED
    return Phase.CLOSED
\end{lstlisting}
\end{frame}

% ------------------------------------------------------------
\begin{frame}[fragile, shrink=10]{Service: registracija s provjerom faze}
\begin{lstlisting}[language=python]
async def create_registration(
    db: AsyncSession,
    comp_id: int,
    body: RegistrationCreate,
    current_user: User,
) -> Registration:
    competition = await competition_repo.get_by_id(db, comp_id)
    if not competition:
        raise AppError("not_found", "Natjecanje ne postoji", 404)

    phase = get_competition_phase(competition)
    if phase != Phase.OPEN:
        raise AppError(
            "deadline_passed",
            "Rok za nove prijave je istekao",
            400,
        )

    lifter = await lifter_repo.get_by_id(db, body.lifter_id)
    if current_user.role == "club":
        if lifter.club_id != current_user.club_id:
            raise AppError("forbidden", "Pristup odbijen", 403)

    reg = Registration(
        competition_id=comp_id, **body.model_dump()
    )
    return await registration_repo.create(db, reg)
\end{lstlisting}
\end{frame}

% ------------------------------------------------------------
\begin{frame}[fragile,shrink=5]{Service: odjava natjecatelja (soft status)}
\begin{lstlisting}[language=python]
async def withdraw_registration(
    db: AsyncSession,
    comp_id: int,
    registration_id: int,
    current_user: User,
) -> Registration:
    reg = await registration_repo.get_by_id(db, registration_id)
    if not reg or reg.competition_id != comp_id:
        raise AppError("not_found", "Prijava ne postoji", 404)
    if reg.status == "withdrawn":
        raise AppError("already_withdrawn", "Vec odjavljen", 400)

    # ownership
    lifter = await lifter_repo.get_by_id(db, reg.lifter_id)
    if current_user.role == "club":
        if lifter.club_id != current_user.club_id:
            raise AppError("forbidden", "Pristup odbijen", 403)

    phase = get_competition_phase(
        await competition_repo.get_by_id(db, comp_id)
    )
    if phase == Phase.CLOSED:
        raise AppError(
            "deadline_passed",
            "Izmjene nisu moguce nakon finalnog roka",
            400,
        )

    reg.status = "withdrawn"
    await db.flush()
    return reg
\end{lstlisting}

\vspace{0.4em}
Soft status: ne brišemo red, mijenjamo \texttt{status} u \texttt{"withdrawn"}. Čuvamo audit trail.
\end{frame}

% ------------------------------------------------------------
\begin{frame}{Matrica operacija po fazama}
\begin{tabular}{@{}lcccc@{}}
  \toprule
  \textbf{Operacija} & \textbf{OPEN} & \textbf{PRELIM} & \textbf{CLOSED} & \textbf{Error} \\
  \midrule
  Nova prijava         & da & ne & ne & 400 \\
  Promjena kategorije  & da & da & ne & 400 \\
  Odjava natjecatelja  & da & da & ne & 400 \\
  Pregled prijava      & da & da & da & - \\
  \bottomrule
\end{tabular}

\vspace{0.6em}
Svaka operacija u serviceu mora provjeriti fazu prije izvršenja.
\end{frame}

% ------------------------------------------------------------
\begin{frame}[fragile,shrink=5]{Promjena kategorije: OPEN i PRELIM\_PASSED}
\begin{lstlisting}[language=python]
async def update_registration_category(
    db: AsyncSession,
    comp_id: int,
    registration_id: int,
    new_category: str,
    current_user: User,
) -> Registration:
    reg = await registration_repo.get_by_id(db, registration_id)
    if not reg or reg.competition_id != comp_id:
        raise AppError("not_found", "Prijava ne postoji", 404)

    # ownership
    lifter = await lifter_repo.get_by_id(db, reg.lifter_id)
    if current_user.role == "club":
        if lifter.club_id != current_user.club_id:
            raise AppError("forbidden", "Pristup odbijen", 403)

    phase = get_competition_phase(
        await competition_repo.get_by_id(db, comp_id)
    )
    if phase == Phase.CLOSED:
        raise AppError(
            "deadline_passed",
            "Izmjene nisu moguce nakon finalnog roka",
            400,
        )

    reg.category = new_category
    await db.flush()
    return reg
\end{lstlisting}
\end{frame}

% ------------------------------------------------------------
\begin{frame}{Hvatanje IntegrityError: dupla prijava}
Unique constraint na \texttt{(lifter\_id, competition\_id)} baca \texttt{IntegrityError} pri duplom pokušaju.

\vspace{0.6em}
Dva pristupa:
\begin{enumerate}
  \item Provjeri u kodu prije inserta (query + if) - ne štiti od race conditiona
  \item Hvati \texttt{IntegrityError} nakon flusha - sigurno, baza je autoritet
\end{enumerate}

\vspace{0.6em}
Koristimo oba: query daje razumljivu grešku, constraint je sigurnosna mreža.
\end{frame}

% ------------------------------------------------------------
\begin{frame}[fragile,shrink=5]{Hvatanje IntegrityError u serviceu}
\begin{lstlisting}[language=python]
from sqlalchemy.exc import IntegrityError

async def create_registration(db, comp_id, body, current_user):
    # ... faza i ownership provjere ...

    reg = Registration(competition_id=comp_id, **body.model_dump())
    try:
        return await registration_repo.create(db, reg)
    except IntegrityError:
        raise AppError(
            "duplicate_registration",
            "Natjecatelj je vec prijavljen na ovo natjecanje",
            409,
        )
\end{lstlisting}

\vspace{0.4em}
409 Conflict - klijent zna da resurs već postoji i može prilagoditi UI.
\end{frame}

% ------------------------------------------------------------
\begin{frame}{Domain errori: dosljedan format}
Svaki error ima tri dijela:

\vspace{0.4em}
\begin{tabular}{@{}llp{6cm}@{}}
  \toprule
  \textbf{Polje} & \textbf{Tip} & \textbf{Primjer} \\
  \midrule
  \texttt{code} & string & \texttt{deadline\_passed} \\
  \texttt{message} & string & \enquote{Rok za nove prijave je istekao} \\
  \texttt{status\_code} & int & 400 \\
  \bottomrule
\end{tabular}

\vspace{0.6em}
\texttt{code} je strojno čitljiv - frontend može reagirati na njega (\texttt{switch} po codeu).
\texttt{message} je za čovjeka.
\end{frame}

% ------------------------------------------------------------
\begin{frame}[fragile,shrink=5]{Exception handler: tko AppError pretvara u HTTP}
Service baca \texttt{AppError}, ali netko ga mora pretvoriti u JSON response. To radi \textbf{globalni exception handler} registriran u \codefile{main.py}:

\begin{lstlisting}[language=python]
# app/main.py
from fastapi.responses import JSONResponse
from app.core.errors import AppError

@app.exception_handler(AppError)
async def app_error_handler(request, exc: AppError):
    return JSONResponse(
        status_code=exc.status_code,
        content={"code": exc.code, "message": exc.message},
    )
\end{lstlisting}

\vspace{0.4em}
Rezultat: nijedan router ne piše try/except za domenske greške. Service baca, handler formatira - svaki error u projektu ima isti oblik. Za Pydanticov 422 FastAPI već ima ugrađeni handler; po želji ga možete zamijeniti da vraća isti \texttt{code}/\texttt{message} format.
\end{frame}

% ------------------------------------------------------------
\begin{frame}{Pregled svih error codeova u projektu}
\begin{tabular}{@{}llll@{}}
  \toprule
  \textbf{Code} & \textbf{HTTP} & \textbf{Kad} \\
  \midrule
  \texttt{invalid\_credentials} & 401 & login s krivom lozinkom \\
  \texttt{token\_expired} & 401 & istekao JWT \\
  \texttt{forbidden} & 403 & nema pristup (rola/ownership) \\
  \texttt{not\_found} & 404 & resurs ne postoji \\
  \texttt{duplicate\_registration} & 409 & dupla prijava \\
  \texttt{already\_withdrawn} & 400 & ponovna odjava vec odjavljenog \\
  \texttt{deadline\_passed} & 400 & operacija izvan dozvoljene faze \\
  \texttt{validation\_error} & 422 & krivi format inputa \\
  \bottomrule
\end{tabular}
\end{frame}

% ------------------------------------------------------------
\begin{frame}{Testiranje poslovnih pravila: strategija}
Tri tipa testova za pravila:
\begin{enumerate}
  \item \textbf{Happy path}: operacija uspije kad su uvjeti zadovoljeni
  \item \textbf{Prekršaj pravila}: operacija je odbijena s točnim error codeom
  \item \textbf{Granični slučaj}: operacija na samoj granici roka
\end{enumerate}

\vspace{0.6em}
Za rokove je ključno kontrolirati \enquote{sada} u testovima - ne želimo da testovi ovise o stvarnom satu.
\end{frame}

% ------------------------------------------------------------
\begin{frame}[fragile,shrink=5]{Freezegun: kontrola vremena u testovima}
\begin{lstlisting}[language=python]
from freezegun import freeze_time

@freeze_time("2026-03-01 10:00:00")
def test_registration_before_prelim(client, admin_token):
    # prelim_deadline je 2026-03-15
    r = client.post(
        "/competitions/1/registrations",
        json={"lifter_id": 1, "category": "83"},
        headers={"Authorization": f"Bearer {admin_token}"},
    )
    assert r.status_code == 201

@freeze_time("2026-03-20 10:00:00")
def test_registration_after_prelim(client, admin_token):
    r = client.post(
        "/competitions/1/registrations",
        json={"lifter_id": 1, "category": "83"},
        headers={"Authorization": f"Bearer {admin_token}"},
    )
    assert r.status_code == 400
    assert r.json()["code"] == "deadline_passed"
\end{lstlisting}
\end{frame}

% ------------------------------------------------------------
\begin{frame}{Alternativa freezegunu: injektiranje vremena}
\begin{itemize}
  \item Umjesto \texttt{datetime.now()} u kodu, primamo \texttt{now} kao parametar
  \item Service metoda: \texttt{def create\_registration(db, body, user, now=None)}
  \item U produkciji: \texttt{now} je \texttt{None} $\rightarrow$ koristi se \texttt{datetime.now(utc)}
  \item U testu: šaljemo fiksnu vrijednost
\end{itemize}

\vspace{0.6em}
Oba pristupa su validna. Freezegun je praktičniji za integration testove, injektiranje za unit testove.
\end{frame}

% ------------------------------------------------------------
\begin{frame}[fragile, shrink=5]{Test: odjava prije i nakon final roka}
\begin{lstlisting}[language=python]
@freeze_time("2026-03-20 10:00:00")
def test_withdraw_between_prelim_and_final(
    client, club_token, registration_id
):
    # final_deadline je 2026-04-01
    r = client.post(
        f"/competitions/1/registrations/{registration_id}/withdraw",
        headers={"Authorization": f"Bearer {club_token}"},
    )
    assert r.status_code == 200
    assert r.json()["status"] == "withdrawn"

@freeze_time("2026-04-05 10:00:00")
def test_withdraw_after_final(
    client, club_token, registration_id
):
    r = client.post(
        f"/competitions/1/registrations/{registration_id}/withdraw",
        headers={"Authorization": f"Bearer {club_token}"},
    )
    assert r.status_code == 400
    assert r.json()["code"] == "deadline_passed"
\end{lstlisting}
\end{frame}

% ------------------------------------------------------------
\begin{frame}[fragile,shrink=5]{Test: dupla prijava daje 409}
\begin{lstlisting}[language=python]
@freeze_time("2026-03-01 10:00:00")
def test_duplicate_registration(client, admin_token):
    payload = {"lifter_id": 1, "category": "83"}
    headers = {"Authorization": f"Bearer {admin_token}"}
    url = "/competitions/1/registrations"

    r1 = client.post(url, json=payload, headers=headers)
    assert r1.status_code == 201

    r2 = client.post(url, json=payload, headers=headers)
    assert r2.status_code == 409
    assert r2.json()["code"] == "duplicate_registration"
\end{lstlisting}
\end{frame}

% ------------------------------------------------------------
\begin{frame}{Što pokrivamo testovima prije obrane}
\begin{enumerate}
  \item Auth: login uspješan, login neuspješan, pristup bez tokena
  \item Role: admin smije, club ne smije (i obratno)
  \item Ownership: club A ne smije pristupiti podacima kluba B
  \item Rokovi: registracija prije/između/nakon rokova
  \item Constraint: dupla prijava daje 409
  \item Validacija: krivi input daje 422
\end{enumerate}

\vspace{0.6em}
Na obrani mogu pokrenuti testove i pitati \textbf{što točno svaki test pokriva}.
\end{frame}

% ------------------------------------------------------------
\begin{frame}{Commitovi za Predavanje 6}
\begin{itemize}
  \item \texttt{feat(api): add Phase enum and get\_competition\_phase helper}
  \item \texttt{feat(api): enforce phase rules on registration create and category update}
  \item \texttt{feat(api): add withdraw endpoint with status transition and phase check}
  \item \texttt{test(api): add phase-based tests with freezegun (OPEN/PRELIM/CLOSED)}
  \item \texttt{test(api): add withdrawal and edge case tests}
\end{itemize}

\vspace{0.4em}
\textbf{Na satu:} prolazimo diffeve, pokrećemo testove, simuliramo scenarije u Swaggeru.
\end{frame}

% ------------------------------------------------------------
\begin{frame}{Rekapitulacija: cijeli backend stack}
\begin{enumerate}
  \item \textbf{P1} - Skeleton, zdravlje, konfiguracija
  \item \textbf{P2} - Modeli, ORM, migracije, seed
  \item \textbf{P3} - Autentikacija (JWT, bcrypt, login/refresh)
  \item \textbf{P4} - Autorizacija (role, ownership, 401/403)
  \item \textbf{P5} - CRUD, scheme, validacija, paginacija
  \item \textbf{P6} - Poslovna pravila, rokovi, workflow
\end{enumerate}

\vspace{0.6em}
Sve zajedno: korisnik se prijavi, ima rolu, vidi samo svoje, poštuje rokove, dobiva jasne greške.
\end{frame}

% ------------------------------------------------------------
\begin{frame}{Priprema za obranu (16.11.)}
\begin{enumerate}
  \item Pokreni projekt: \texttt{docker compose up -d db} + \texttt{uvicorn}
  \item Pokaži CRUD u Swaggeru
  \item Pokreni \texttt{pytest} - svi testovi prolaze
  \item Budi spreman objasniti:
  \begin{itemize}
    \item Gdje provjeravaš ownership?
    \item Kako backend zabranjuje izmjene nakon roka?
    \item Zašto je status code baš taj (401/403/409)?
    \item Što test pokriva i zašto je bitan?
  \end{itemize}
\end{enumerate}
\end{frame}

% ------------------------------------------------------------
\begin{frame}{Mini promjena na obrani}
Na obrani mogu tražiti malu promjenu uživo:
\begin{itemize}
  \item Dodaj novo polje u model + migracija
  \item Dodaj constraint ili provjeru u service
  \item Promijeni error poruku ili status code
  \item Dodaj novi test za edge case
\end{itemize}

\vspace{0.6em}
Nema potrebe za panikom. Ako razumijete slojeve i znate gdje što živi, promjena traje par minuta.
\end{frame}

% ------------------------------------------------------------
\begin{frame}{Što ide sljedeće (Predavanje 7)}
\begin{itemize}
  \item Arhitektura frontend aplikacija
  \item Vue 3: komponente, reaktivnost, Composition API
  \item Vite dev server i build pipeline
  \item Struktura SPA projekta
\end{itemize}

\vspace{0.8em}
\begin{center}
\Large Pitanja?
\end{center}
\end{frame}

\end{document}
